% questions.tex —— 2013 年·湖北卷·文科 · 题目内容片段（无答案）
% 由 tools/generate_projects.py 依来源总账自 2013/2013_数学真题_湖北文_试题.md 生成，勿手改（改源 MD 后重跑）

\gkpart{一、选择题}

\begin{question}
\gkti{1}已知全集 $U=\{1,2,3,4,5\}$，集合 $A=\{1,2\}$，$B=\{2,3,4\}$，则 $B\cap\complement_U A=$\mbox{（　　）}

（A）$\{2\}$；（B）$\{3,4\}$；（C）$\{1,4,5\}$；（D）$\{2,3,4,5\}$。

\end{question}
\begin{question}
\gkti{2}已知 $0<\theta<\dfrac\pi4$，则双曲线 $C_1:\dfrac{x^2}{\sin^2\theta}-\dfrac{y^2}{\cos^2\theta}=1$ 与 $C_2:\dfrac{y^2}{\cos^2\theta}-\dfrac{x^2}{\sin^2\theta}=1$ 的\mbox{（　　）}

（A）实轴长相等；（B）虚轴长相等；（C）离心率相等；（D）焦距相等。

\end{question}
\begin{question}
\gkti{3}在一次跳伞训练中，甲、乙两位学员各跳一次。设命题 $p$ 是“甲降落在指定范围”，$q$ 是“乙降落在指定范围”，则命题“至少有一位学员没有降落在指定范围”可表示为\mbox{（　　）}

（A）$(\neg p)\lor(\neg q)$；（B）$p\lor(\neg q)$；（C）$(\neg p)\land(\neg q)$；（D）$p\lor q$。

\end{question}
\begin{question}
\gkti{4}四名同学根据各自的样本数据研究变量 $x$，$y$ 之间的相关关系，并求得回归直线方程，分别得到以下四个结论：

① $y$ 与 $x$ 负相关且 $\hat y=2.347x-6.423$；

② $y$ 与 $x$ 负相关且 $\hat y=-3.476x+5.648$；

③ $y$ 与 $x$ 正相关且 $\hat y=5.437x+8.493$；

④ $y$ 与 $x$ 正相关且 $\hat y=-4.326x-4.578$。

其中一定不正确的结论的序号是\mbox{（　　）}

（A）①②；（B）②③；（C）③④；（D）①④。

\Needspace{0.4\textheight}
\end{question}
\begin{question}
\gkti{5}小明骑车上学，开始时匀速行驶，途中因交通堵塞停留了一段时间后，为了赶时间加快速度行驶。与以上事件吻合得最好的图象是\mbox{（　　）}
\nopagebreak[4]
\begin{center}% 2013 湖北卷文科第 5 题：四个“距学校的距离—时间”图象。
\begin{tikzpicture}[font=\footnotesize, >=Stealth, line width=0.5pt]
  % (A)：距离从 0 开始增大，先匀速，后短暂平缓，再加速上升
  \begin{scope}[shift={(0,0)}]
    \draw[->] (0,0) -- (0,2.05) node[above] {距学校的距离};
    \draw[->] (0,0) -- (3.6,0) node[below right] {时间};
    \node[below left] at (0,0) {$O$};
    \draw[thick] (0,0) -- (1.25,0.58) -- (2.05,0.66)
      .. controls (2.45,0.68) and (2.70,0.95) .. (2.98,1.58);
    \node at (1.8,-0.55) {（A）};
  \end{scope}
  % (B)：先较快下降，短暂变缓，再匀速下降
  \begin{scope}[shift={(5.75,0)}]
    \draw[->] (0,0) -- (0,2.05) node[above] {距学校的距离};
    \draw[->] (0,0) -- (3.6,0) node[below right] {时间};
    \node[below left] at (0,0) {$O$};
    \draw[thick] (0,1.58) -- (1.55,1.08) -- (2.10,1.02) -- (3.25,0.16);
    \node at (1.8,-0.55) {（B）};
  \end{scope}
  % (C)：匀速下降，中途停留（水平），再加速下降
  \begin{scope}[shift={(0,-3.75)}]
    \draw[->] (0,0) -- (0,2.05) node[above] {距学校的距离};
    \draw[->] (0,0) -- (3.6,0) node[below right] {时间};
    \node[below left] at (0,0) {$O$};
    \draw[thick] (0,1.60) -- (1.20,1.10) -- (2.05,1.08)
      .. controls (2.42,1.06) and (2.62,0.72) .. (2.95,0.05);
    \node at (1.8,-0.55) {（C）};
  \end{scope}
  % (D)：匀速下降，随后直接加速下降（无停留）
  \begin{scope}[shift={(5.75,-3.75)}]
    \draw[->] (0,0) -- (0,2.05) node[above] {距学校的距离};
    \draw[->] (0,0) -- (3.6,0) node[below right] {时间};
    \node[below left] at (0,0) {$O$};
    \draw[thick] (0,1.60) -- (1.85,0.92) -- (2.85,0.05);
    \node at (1.8,-0.55) {（D）};
  \end{scope}
\end{tikzpicture}
\end{center}
\nopagebreak[4]
\end{question}
\begin{question}
\gkti{6}将函数 $y=\sqrt3\cos x+\sin x$（$x\in\mathbf R$）的图象向左平移 $m$（$m>0$）个单位长度后，所得到的图象关于 $y$ 轴对称，则 $m$ 的最小值是\mbox{（　　）}

（A）$\dfrac\pi{12}$；（B）$\dfrac\pi6$；（C）$\dfrac\pi3$；（D）$\dfrac{5\pi}6$。

\end{question}
\begin{question}
\gkti{7}已知点 $A(-1,1)$，$B(1,2)$，$C(-2,-1)$，$D(3,4)$，则向量 $\overrightarrow{AB}$ 在 $\overrightarrow{CD}$ 方向上的投影为\mbox{（　　）}

（A）$\dfrac{3\sqrt2}2$；（B）$\dfrac{3\sqrt{15}}2$；（C）$-\dfrac{3\sqrt2}2$；（D）$-\dfrac{3\sqrt{15}}2$。

\end{question}
\begin{question}
\gkti{8}$x$ 为实数，$[x]$ 表示不超过 $x$ 的最大整数，则函数 $f(x)=x-[x]$ 在 $\mathbf R$ 上为\mbox{（　　）}

（A）奇函数；（B）偶函数；（C）增函数；（D）周期函数。

\end{question}
\begin{question}
\gkti{9}某旅行社租用 $A$，$B$ 两种型号的客车安排 $900$ 名客人旅行，$A$，$B$ 两种车辆的载客量分别为 $36$ 人和 $60$ 人，租金分别为 $1600$ 元/辆和 $2400$ 元/辆，旅行社要求租车总数不超过 $21$ 辆，且 $B$ 型车不多于 $A$ 型车 $7$ 辆，则租金最少为\mbox{（　　）}

（A）$31200$ 元；（B）$36000$ 元；（C）$36800$ 元；（D）$38400$ 元。

\end{question}
\begin{question}
\gkti{10}已知函数 $f(x)=x(\ln x-ax)$ 有两个极值点，则实数 $a$ 的取值范围是\mbox{（　　）}

（A）$(-\infty,0)$；（B）$\left(0,\dfrac12\right)$；（C）$(0,1)$；（D）$(0,+\infty)$。

\end{question}
\gkpart{二、填空题}

\begin{question}
\gkti{11}$\mathrm i$ 为虚数单位，设复数 $z_1$，$z_2$ 在复平面内对应的点关于原点对称，若 $z_1=2-3\mathrm i$，则 $z_2=$\gkblank{}。

\end{question}
\begin{question}
\gkti{12}某学员在一次射击测试中射靶 $10$ 次，命中环数如下：

\nopagebreak[4]
\[7,\ 8,\ 7,\ 9,\ 5,\ 4,\ 9,\ 10,\ 7,\ 4\]

则：

（1） 平均命中环数为\gkblank{}；

（2） 命中环数的标准差为\gkblank{}。

\Needspace{0.4\textheight}
\end{question}
\begin{question}
\gkti{13}阅读如图所示的程序框图，运行相应的程序。若输入 $m$ 的值为 $2$，则输出的结果 $i=$\gkblank{}。
\nopagebreak[4]
\begin{center}% 2013 湖北卷文科第 13 题：程序框图。
\begin{tikzpicture}[font=\footnotesize, >=Stealth, line width=0.45pt,
  node distance=4.2mm,
  box/.style={draw, minimum height=1.7em, inner xsep=5pt},
  term/.style={draw, rounded corners=5pt, minimum height=1.7em, inner xsep=8pt},
  io/.style={draw, trapezium, trapezium left angle=70, trapezium right angle=110,
             minimum height=1.7em, inner xsep=5pt},
  dec/.style={draw, diamond, aspect=2.1, inner sep=1pt}]
  \node[term] (start) {开始};
  \node[io, below=5mm of start] (input) {输入 $m$};
  \node[box, below=5mm of input] (init) {$A=1$，$B=1$，$i=0$};
  \node[box, below=6mm of init] (inc) {$i=i+1$};
  \node[box, below=4.5mm of inc] (am) {$A=A\times m$};
  \node[box, below=4.5mm of am] (bi) {$B=B\times i$};
  \node[dec, below=6mm of bi] (cond) {$A<B?$};
  \node[io, below=7mm of cond] (output) {输出 $i$};
  \node[term, below=5mm of output] (end) {结束};
  \draw[->] (start) -- (input);
  \draw[->] (input) -- (init);
  \draw[->] (init) -- (inc);
  \draw[->] (inc) -- (am);
  \draw[->] (am) -- (bi);
  \draw[->] (bi) -- (cond);
  \draw[->] (cond) -- node[right] {是} (output);
  \draw[->] (output) -- (end);
  % “否”分支：从判断框右侧引出，向上回到 i=i+1 之前的入口
  \draw[->] (cond.east) -- node[above] {否} ++(1.15,0)
    |- ($(inc.north)!0.5!(init.south)$);
\end{tikzpicture}
\end{center}
\nopagebreak[4]
\end{question}
\begin{question}
\gkti{14}已知圆 $O:x^2+y^2=5$，直线 $l:x\cos\theta+y\sin\theta=1\left(0<\theta<\dfrac\pi2\right)$。设圆 $O$ 上到直线 $l$ 的距离等于 $1$ 的点的个数为 $k$，则 $k=$\gkblank{}。

\end{question}
\begin{question}
\gkti{15}在区间 $[-2,4]$ 上随机地取一个数 $x$，若 $x$ 满足 $|x|\leqslant m$ 的概率为 $\dfrac56$，则 $m=$\gkblank{}。

\end{question}
\begin{question}
\gkti{16}我国古代数学名著《数书九章》中有“天池盆测雨”题：在下雨时，用一个圆台形的天池盆接雨水，天池盆盆口直径为二尺八寸，盆底直径为一尺二寸，盆深一尺八寸，若盆中积水深九寸，则平地降雨量是\gkblank{}寸。

（注：① 平地降雨量等于盆中积水体积除以盆口面积；② 一尺等于十寸）

\Needspace{0.4\textheight}
\end{question}
\begin{question}
\gkti{17}在平面直角坐标系中，若点 $P(x,y)$ 的坐标 $x$，$y$ 均为整数，则称点 $P$ 为格点，若一个多边形的顶点全是格点，则称该多边形为格点多边形，格点多边形的面积为 $S$，其内部的格点数记为 $N$，边界上的格点数记为 $L$。例如图中 $\triangle ABC$ 是格点三角形，对应的 $S=1$，$N=0$，$L=4$。
\nopagebreak[4]
\begin{center}% 2013 湖北卷文科第 17 题：格点四边形 DEFG 与格点三角形 ABC。
\begin{tikzpicture}[font=\footnotesize, >=Stealth, line width=0.45pt, scale=0.9]
  % 网格
  \draw[gray!45, very thin] (0,0) grid (6,6);
  % 坐标轴
  \draw[->] (0,0) -- (6.55,0) node[below right] {$x$};
  \draw[->] (0,0) -- (0,6.55) node[left] {$y$};
  \node[below left] at (0,0) {$O$};
  \foreach \x in {1,...,5} {
    \draw (\x,0) -- (\x,-0.06);
    \node[below] at (\x,-0.09) {$\x$};
  }
  \foreach \y in {1,...,5} {
    \draw (0,\y) -- (-0.06,\y);
    \node[left] at (-0.09,\y) {$\y$};
  }
  % 格点三角形 ABC
  \draw[thick] (1,2) -- (1,1) -- (3,1) -- cycle;
  \fill (1,2) circle (1.6pt) (1,1) circle (1.6pt) (3,1) circle (1.6pt);
  \node[above left] at (1,2) {$B$};
  \node[below left] at (1,1) {$C$};
  \node[below] at (3,1) {$A$};
  % 格点四边形 DEFG
  \draw[thick] (2,3) -- (4,3) -- (2,5) -- (1,4) -- cycle;
  \fill (2,3) circle (1.6pt) (4,3) circle (1.6pt) (2,5) circle (1.6pt) (1,4) circle (1.6pt);
  \node[below left] at (2,3) {$D$};
  \node[right] at (4,3) {$E$};
  \node[above] at (2,5) {$F$};
  \node[left] at (1,4) {$G$};
  % 图中的其余格点
  \fill (2,4) circle (1.3pt) (3,4) circle (1.3pt) (3,3) circle (1.3pt);
\end{tikzpicture}
\end{center}
\nopagebreak[4]
（1） 图中格点四边形 $DEFG$ 对应的 $S$，$N$，$L$ 分别是\gkblank{}；

（2） 已知格点多边形的面积可表示为 $S=aN+bL+c$，其中 $a$，$b$，$c$ 为常数。若某格点多边形对应的 $N=71$，$L=18$，则 $S=$\gkblank{}。（用数值作答）

\end{question}
\gkpart{三、解答题}

\begin{question}
\gkti{18}在 $\triangle ABC$ 中，角 $A$，$B$，$C$ 对应的边分别是 $a$，$b$，$c$，已知 $\cos2A-3\cos(B+C)=1$。

（1） 求角 $A$ 的大小；

（2） 若 $\triangle ABC$ 的面积 $S=5\sqrt3$，$b=5$，求 $\sin B\sin C$ 的值。

\end{question}
\begin{question}
\gkti{19}已知 $S_n$ 是等比数列 $\{a_n\}$ 的前 $n$ 项和，$S_4$，$S_2$，$S_3$ 成等差数列，且 $a_2+a_3+a_4=-18$。

（1） 求数列 $\{a_n\}$ 的通项公式；

（2） 是否存在正整数 $n$，使得 $S_n\geqslant2013$？若存在，求出符合条件的所有的 $n$ 的集合；若不存在，说明理由。

\Needspace{0.4\textheight}
\end{question}
\begin{question}
\gkti{20}如图，某地质队自水平地面 $A$，$B$，$C$ 三处垂直向地下钻探，自 $A$ 点向下钻到 $A_1$ 处发现矿藏，再继续下钻到 $A_2$ 处后下面已无矿，从而得到在 $A$ 处正下方的矿层厚度为 $A_1A_2=d_1$，同样可得在 $B$，$C$ 处正下方的矿层厚度分别为 $B_1B_2=d_2$，$C_1C_2=d_3$，且 $d_1<d_2<d_3$，过 $AB$，$AC$ 的中点 $M$，$N$ 且与直线 $AA_2$ 平行的平面截多面体 $A_1B_1C_1-A_2B_2C_2$ 所得的截面 $DEFG$ 为该多面体的一个中截面，其面积记为 $S_{\text{\text{中}}}$。
\nopagebreak[4]
\begin{center}% 2013 湖北卷文科第 20 题：矿层多面体 A1B1C1-A2B2C2 与中截面 DEFG。
\begin{tikzpicture}[font=\footnotesize, >=Stealth, line width=0.5pt, scale=0.9]
  % 地面三角形 ABC（在水平面内）
  \coordinate (A) at (0,0);
  \coordinate (B) at (6,0);
  \coordinate (C) at (4.0,-1.30);
  % 矿层顶面三角形 A1B1C1
  \coordinate (A1) at (0,-1.70);
  \coordinate (B1) at (6,-1.90);
  \coordinate (C1) at (4.0,-3.00);
  % 矿层底面三角形 A2B2C2
  \coordinate (A2) at (0,-2.60);
  \coordinate (B2) at (6,-3.50);
  \coordinate (C2) at (4.0,-5.00);
  % 中点 M, N 与中截面顶点
  \coordinate (M) at (3,0);
  \coordinate (N) at (2,-0.65);
  \coordinate (D) at (3,-1.80);
  \coordinate (E) at (3,-3.05);
  \coordinate (G) at (2,-2.35);
  \coordinate (F) at (2,-3.80);

  % 地面与钻探竖线
  \draw (A) -- (B) -- (C) -- cycle;
  \draw (A) -- (A1) -- (A2);
  \draw (B) -- (B1) -- (B2);
  \draw (C) -- (C1) -- (C2);
  % 顶面、底面的边（背面用虚线）
  \draw[dashed] (A1) -- (B1);
  \draw (B1) -- (C1);
  \draw[dashed] (C1) -- (A1);
  \draw[dashed] (A2) -- (B2);
  \draw (B2) -- (C2);
  \draw[dashed] (C2) -- (A2);

  % 中截面 DEFG
  \fill[gray!18] (D) -- (E) -- (F) -- (G) -- cycle;
  \draw[very thick] (D) -- (E) -- (F) -- (G) -- cycle;
  % 截面所在平面与 M, N 的连线（辅助虚线）
  \draw[dashed] (M) -- (D);
  \draw[dashed] (N) -- (G);
  \draw[dashed] (M) -- (N);

  % 标注
  \node[above left] at (A) {$A$};
  \node[above right] at (B) {$B$};
  \node[right] at (C) {$C$};
  \node[left] at (A1) {$A_1$};
  \node[right] at (B1) {$B_1$};
  \node[right] at (C1) {$C_1$};
  \node[left] at (A2) {$A_2$};
  \node[right] at (B2) {$B_2$};
  \node[below] at (C2) {$C_2$};
  \node[above] at (M) {$M$};
  \node[left] at (N) {$N$};
  \node[right] at ($(D)+(0.30,-0.26)$) {$D$};
  \node[right] at ($(E)+(0.22,-0.10)$) {$E$};
  \node[left] at (F) {$F$};
  \node[left] at (G) {$G$};
\end{tikzpicture}
\end{center}
\nopagebreak[4]
（1） 证明：中截面 $DEFG$ 是梯形；

（2） 在 $\triangle ABC$ 中，记 $BC=a$，$BC$ 边上的高为 $h$，面积为 $S$。在估测三角形 $ABC$ 区域内正下方的矿藏储量（即多面体 $A_1B_1C_1-A_2B_2C_2$ 的体积 $V$）时，可用近似公式 $V_{\text{\text{估}}}=S_{\text{\text{中}}}\cdot h$ 来估算。已知 $V=\dfrac13(d_1+d_2+d_3)S$，试判断 $V_{\text{\text{估}}}$ 与 $V$ 的大小关系，并加以证明。

\end{question}
\begin{question}
\gkti{21}设 $a>0$，$b>0$，已知函数 $f(x)=\dfrac{ax+b}{x+1}$。

（1） 当 $a\neq b$ 时，讨论函数 $f(x)$ 的单调性；

（2） 当 $x>0$ 时，称 $f(x)$ 为 $a$，$b$ 关于 $x$ 的加权平均数。

① 判断 $f(1)$，$f\left(\sqrt{\dfrac ba}\right)$，$f\left(\dfrac ba\right)$ 是否成等比数列，并证明 $f\left(\dfrac ba\right)\leqslant f\left(\sqrt{\dfrac ba}\right)$；

② $a$，$b$ 的几何平均数记为 $G$，称 $\dfrac{2ab}{a+b}$ 为 $a$，$b$ 的调和平均数，记为 $H$，若 $H\leqslant f(x)\leqslant G$，求 $x$ 的取值范围。

\Needspace{0.4\textheight}
\end{question}
\begin{question}
\gkti{22}如图，已知椭圆 $C_1$ 与 $C_2$ 的中心在坐标原点 $O$，长轴均为 $MN$ 且在 $x$ 轴上，短轴长分别为 $2m$，$2n$（$m>n$），过原点且不与 $x$ 轴重合的直线 $l$ 与 $C_1$，$C_2$ 的四个交点按纵坐标从大到小依次为 $A$，$B$，$C$，$D$，记 $\lambda=\dfrac mn$，$\triangle BDM$ 和 $\triangle ABN$ 的面积分别为 $S_1$ 和 $S_2$。
\nopagebreak[4]
\begin{center}% 2013 湖北卷文科第 22 题：共焦轴的两椭圆与过原点的直线 l。
\begin{tikzpicture}[font=\footnotesize, >=Stealth, line width=0.5pt, scale=1.0]
  \def\a{4}
  \def\m{2.4}
  \def\n{1.4}
  % 椭圆
  \draw[thick] (0,0) ellipse ({\a} and {\m});
  \draw[thick] (0,0) ellipse ({\a} and {\n});
  % 坐标轴
  \draw[->] (-4.8,0) -- (4.9,0) node[below right] {$x$};
  \draw[->] (0,-3.0) -- (0,3.1) node[left] {$y$};
  \node[below left] at (0,0) {$O$};
  % 四个交点（直线 y=0.75x）
  \coordinate (A) at (2.499,1.874);
  \coordinate (B) at (1.692,1.269);
  \coordinate (C) at (-1.692,-1.269);
  \coordinate (D) at (-2.499,-1.874);
  \coordinate (M) at (-\a,0);
  \coordinate (N) at (\a,0);
  % 直线 l
  \draw[thick] (-4.7,-3.525) -- (4.7,3.525);
  % 两个三角形 BDM 与 ABN 的边
  \draw (M) -- (B);
  \draw (B) -- (N);
  \draw (A) -- (N);
  \draw (M) -- (D);
  \draw (A) -- (B);
  % 交点与端点
  \fill (A) circle (1.6pt) (B) circle (1.6pt) (C) circle (1.6pt) (D) circle (1.6pt);
  \fill (M) circle (1.6pt) (N) circle (1.6pt);
  \node[fill=white, inner sep=1pt, above left] at (A) {$A$};
  \node[fill=white, inner sep=1pt, above left] at (B) {$B$};
  \node[fill=white, inner sep=1pt, above left] at (C) {$C$};
  \node[fill=white, inner sep=1pt, below right] at (D) {$D$};
  \node[fill=white, inner sep=1pt, above left] at (M) {$M$};
  \node[fill=white, inner sep=1pt, above right] at (N) {$N$};
\end{tikzpicture}
\end{center}
\nopagebreak[4]
（1） 当直线 $l$ 与 $y$ 轴重合时，若 $S_1=\lambda S_2$，求 $\lambda$ 的值；

（2） 当 $\lambda$ 变化时，是否存在与坐标轴不重合的直线 $l$，使得 $S_1=\lambda S_2$？并说明理由。

\end{question}